% v6 JHEP-formal companion paper (ECI framework, see paper/eci.tex for v5
% phenomenological companion). Title selected: T3.
%
% Class: jheppub.cls unavailable locally at build time (SISSA 404 on
% 2026-04-22); compiled with RevTeX 4-2 for the preprint stage. Swap to
% jheppub before JHEP submission — see v6/README.md.

\documentclass[aps,prd,superscriptaddress,nofootinbib,floatfix,10pt]{revtex4-2}

\usepackage{amsmath,amssymb,amsthm}
\usepackage{mathtools}
\usepackage{hyperref}
\usepackage{xcolor}
\usepackage{enumitem}
\usepackage{tcolorbox}
\usepackage{tikz}
\usepackage{tikz-cd}

\newtheorem{theorem}{Theorem}
\newtheorem{proposition}{Proposition}
\newtheorem{lemma}{Lemma}
\newtheorem{definition}{Definition}
\newtheorem{assumption}{Assumption}
\theoremstyle{remark}
\newtheorem{motivation}{Motivation}

\newcommand{\Sgen}{S_{\mathrm{gen}}}
\newcommand{\tauR}{\tau_{R}}
\newcommand{\kR}{\kappa_{R}}
\newcommand{\Ck}{\mathcal{C}_{k}}
\newcommand{\PHk}{\mathrm{PH}_{k}}
\newcommand{\dn}{\delta n}
\newcommand{\AR}{\mathcal{A}_{R}}

\begin{document}

\title{From Faulkner--Speranza to a complexity-bounded\\
generalised second law in type-II crossed-product algebras}

\author{Kevin Remondi\`ere}
\email{kevin.remondiere@gmail.com}
\affiliation{Independent Researcher, ORCID 0009-0008-2443-7166}

\date{\today}

\begin{abstract}
We prove an upper bound on the rate of generalised entropy along the
modular (Connes--Tomita--Takesaki) flow $\tauR$ of a type-II
crossed-product observer algebra $\AR$, extending the
Eling--Guedens--Jacobson local Clausius schema
$dS = \delta Q / T + d_{i} S$ to the setting of
Chandrasekaran--Longo--Penington--Witten and
De~Vuyst--Eccles--Held--Kirklin. The bound reads
$\, d\Sgen[R]/d\tauR \le \kR \cdot \Ck[\rho_{R}(\tauR)]
 \cdot \Theta(\PHk[\dn(\tauR)])\,$
where $\Ck$ is the $k$-design complexity of the reduced state,
$\PHk$ is the order-$k$ persistent-homology Betti number of a
dequantised density field $\dn$, and $\Theta$ is a chameleon-type
topological activator. We further establish a tightened logistic
envelope $\, d\Sgen/d\tauR \le \kR \Ck (1 - \Ck/\Ck^{\max})\Theta\,$
in the scrambling regime with
$\Ck^{\max}\sim\exp(S_{R})$ (Prop.~\ref{prop:logistic}), which
dovetails with the Fan (2022) logarithmic Krylov law at full
saturation. The dequantisation map
$\dn(x,\tauR) \equiv \mathrm{Tr}_{R}[\rho_{R}(\tauR)\hat n(x)] - \langle n\rangle$
bridges the type-II factor to classical persistent homology and is a
second formal contribution. The main bound recovers Wall's
monotonicity inequality and the Faulkner--Speranza modular form in
the $\Theta\to 1$ limit. We make no cosmological claim;
phenomenological consequences are deferred to a companion paper.
\end{abstract}

\maketitle

% =============================================================
\section{Introduction}
\label{sec:intro}
% =============================================================

The local Clausius derivation of Einstein's equations by
Jacobson~\cite{Jacobson1995} and its non-equilibrium extension by
Eling--Guedens--Jacobson (EGJ, 2006) cast gravity as thermodynamics on
local causal horizons, with an internal-production term
$d_{i}S$ that encodes the departure from equilibrium. Two recent
developments invite a modular-time refinement of that schema. First,
the type-II crossed-product algebras $\AR$ of
CLPW~\cite{CLPW2023} and
DEHK~\cite{DEHK2025a,DEHK2025b} promote observer-dependent generalised
entropy from a formal pairing to a von Neumann trace. Second, the
differential GSL of Faulkner--Speranza~\cite{FaulknerSperanza2024} and
Kirklin~\cite{Kirklin2025} upgrades Wall's monotonicity to a modular
inequality for relative entropy.

The present paper combines these two lines. We propose a
differential upper bound on $d\Sgen/d\tauR$ whose source is the
$k$-design complexity $\Ck[\rho_{R}]$ of the reduced state and whose
topological activator is an exponential function of the persistent-
homology Betti number $\PHk[\dn]$ of a dequantised density field $\dn$.
The result (Eq.~\eqref{eq:main}) is a formal statement; no
cosmological prediction is claimed. Phenomenological falsifiers and
observational forecasts are developed in the phenomenological
companion paper~\cite{Remondiere2026ECIv5}.

We present two formal results on equal footing:
(i)~a differential-inequality upper bound on $d\Sgen/d\tauR$ along the
modular flow of $\AR$, whose source is $k$-design
complexity~\cite{MaHuang2025,Haferkamp2022} (replacing geometric shear
squared of EGJ and logarithmic Krylov complexity of Fan 2022; the
logarithmic form is recovered at saturation, Sec.~\ref{sec:relation});
(ii)~an explicit dequantisation map from the type-II factor $\AR$ to a
classical scalar field $\dn$, which resolves the category mismatch
between quantum reduced states and classical persistent-homology
functors and is an independently usable construction. A minor but
distinct structural novelty is that the bound is modular-time
differential, versus the static type-II second law of CLPW/DEHK.

% =============================================================
\section{Setup}
\label{sec:setup}
% =============================================================

Let $\AR$ denote the type-II crossed-product algebra associated with an
observer-with-clock in region $R$, as constructed by
CLPW~\cite{CLPW2023} and extended by
DEHK~\cite{DEHK2025a,DEHK2025b}. The Connes cocycle fixes, up to an
additive constant, a faithful normal semifinite trace, and we denote by
$\sigma^{R}_{\tauR}$ the one-parameter automorphism of $\AR$ induced by
Tomita--Takesaki modular flow; $\tauR$ is the thermal (Connes--Rovelli)
time for a thermofield double state on $\AR$. The reduced state
$\rho_{R}(\tauR)\equiv\sigma^{R}_{\tauR}(\rho_{R})$ evolves
covariantly.

We adopt the modular-temperature identification
\begin{equation}
 \kR \;\equiv\; 2\pi\,T_{R},
 \label{eq:kappa}
\end{equation}
with $T_{R}$ the local Tomita--Takesaki temperature, so that $\kR$ has
units of nats per modular e-fold and reduces to the Hubble rate $H$ at
the de~Sitter Gibbons--Hawking limit.

The generalised entropy is the type-II functor
$\Sgen[R] = A[R]/(4G_{N}) + S_{\mathrm{matter}}[\rho_{R}]$, well-defined
because the crossed-product factor is semifinite.
The $k$-design complexity $\Ck[\rho_{R}]$ is defined in the
Ma--Huang pseudorandom-unitary (PRU) sense~\cite{MaHuang2025};
it is a dimensionless, monotone, state-dependent functional that
interpolates between circuit complexity~\cite{Haferkamp2022} and
information-theoretic complexity~\cite{CryptoCensorship}.
Operationally, we identify $\Ck[\rho_{R}]$ with the spread complexity of
Caputa--Mag\'an--Patramanis--Tonni 2024~\cite{Caputa2024} evaluated
under the modular flow $e^{-i K_{R}\tau}$ and truncated at the
$k$-design saturation scale; this identification is a sub-postulate
(M1$'$) of M1 below. Concretely, $\Ck[\rho_{R}]$ is the spread complexity
$C_{\mathrm{spread}}(\tauR)$ of~\cite{Caputa2024} computed on the tracial
vacuum of $\AR$ under $e^{-i K_{R}\tauR}$ and truncated at the Haferkamp
linear-growth scale $\tau_{\mathrm{sat}}\sim\exp(S_{R})$.
The spread-complexity identification canonically fixes the gate base up
to the polynomial equivalence class of right-invariant complexity
geometries~\cite{Brown2024Polynomial}; inequality~\eqref{eq:main} is
preserved under any rescaling within that class.
M1 admits an equivalent reformulation as a capacity-limited rate law
in a Gibbs complexity ensemble: with
$\hat{R}_{n} = (\mathcal{C}_{k}^{\max}-n)/\mathcal{C}_{k}^{\max}$ one
has $\langle\hat{R}\rangle = 1-\langle\mathcal{C}_k\rangle/\mathcal{C}_k^{\max}$
exactly, and the logistic envelope~\eqref{eq:logistic} is recovered
under the rate postulate
$d\langle\mathcal{C}_k\rangle/d\tauR \le \kR\langle\mathcal{C}_k\rangle\langle\hat{R}\rangle$;
this framing trades a Brown--Susskind analogy for a maximum-entropy
capacity bound without changing the postulate status of M1.
The Betti number $\PHk[\dn]$ is the order-$k$ persistent-homology count
on a classical density field $\dn$ (defined in Sec.~\ref{sec:dequant}),
following the Yip~et~al.\ cosmological pipeline~\cite{Yip2024}.
The overall structure of the setup, including the dequantisation
branch, is summarised in Fig.~\ref{fig:setup}.

\begin{figure}[htbp]
\centering
\begin{tikzcd}[row sep=large, column sep=large]
\mathrm{QRF}\ R \arrow[r] \arrow[dr, swap]
  & \AR \arrow[r, "\tauR"]
  & \dfrac{d\Sgen}{d\tauR} \\
  & \rho_{R} \arrow[r] \arrow[u, dashed]
  & \dn \arrow[r, "\PHk"]
  & \Theta \arrow[ul, dashed, bend right=15]
\end{tikzcd}
\caption{Setup diagram: a QRF observer $R$ defines the crossed-product
algebra $\AR$, whose modular flow $\tauR$ generates $d\Sgen/d\tauR$.
The dequantisation map bridges $\rho_{R}$ to a classical density field
$\dn$ on which persistent homology $\PHk$ is computed and activator
$\Theta$ evaluated.}
\label{fig:setup}
\end{figure}

% =============================================================
\section{Main inequality}
\label{sec:main}
% =============================================================

\begin{tcolorbox}[colback=blue!3,colframe=blue!50,title=Main result (Theorem~\ref{thm:main})]
\begin{equation}
 \boxed{\;
  \frac{d\Sgen[R]}{d\tauR}
  \;\le\;
  \kR \,\cdot\, \Ck\!\left[\rho_{R}(\tauR)\right]
       \,\cdot\, \Theta\!\left(\PHk[\dn(\tauR)]\right)
 \;}
 \label{eq:main}
\end{equation}
\end{tcolorbox}

By construction $0 \le \Theta(\PHk) \le 1$,\footnote{Note on notation:
Bianconi~2025~\cite{Bianconi2025} uses an auxiliary Lagrange-multiplier
field in the relative-entropy action $S(g\|G)$; our $\Theta$ is a
distinct chameleon-like activator on the topological space of density
fields. The two objects are mathematically and physically unrelated.}
so the RHS is a non-negative upper bound on the non-negative LHS (GSL);
the inequality never reverses sign.

Equation~\eqref{eq:main} is an inequality, not an equality; the
equality attempt was ruled out by an internal multi-agent derivation
audit (three independent routes converged on \emph{ansatz} status;
cf.\ Appendix~\ref{app:computation}). The inequality is tight: it is
saturated in the scrambling regime where
$\Ck\to\exp(S_{R})$, and in that limit recovers the Fan~(2022)
logarithmic Krylov law as the \emph{saturating}, not \emph{derived},
rate (Sec.~\ref{sec:relation}).

The bound rests on three labelled postulates.

\begin{assumption}[M1 --- Modular complexity ansatz, POSTULATE]
\label{ass:M1}
Under modular flow on $\AR$, the internal-production part of
$d\Sgen/d\tauR$ is bounded above by $\kR\,\Ck[\rho_{R}]$. This is a
Brown--Susskind-style ``second law of complexity'' ansatz adapted to
type-II modular time, in the scaffolding proposed by
Caputa--Mag\'an (2024); no type-II theorem is claimed.
\end{assumption}

\begin{assumption}[M2 --- Semi-classical recovery of $\dn$, ANSATZ]
\label{ass:M2}
In the large-$N$ QRF coarse-graining limit, the dequantisation map
(Sec.~\ref{sec:dequant}) produces a mean-zero Gaussian random field
$\dn$ by a central-limit / law-of-large-numbers argument on the
unobserved QRF degrees of freedom. The Gaussian recovery is verified
partially on a toy type-II$_{1}$ factor
(cf.\ \texttt{V6-dequantisation-map.py}); a full proof is outside the
scope of this paper.
\end{assumption}

\begin{assumption}[M3 --- Topological activator profile, CONJECTURAL]
\label{ass:M3}
The activator $\Theta$ takes the chameleon form
$\Theta(\PHk) = \exp\!\left[-(\PHk/\PHk^{\mathrm{c}})^{\alpha}\right]$
with a single power $\alpha \in (0, 0.1]$, inspired by the
Barrow--$\Delta \lesssim 0.1$ fractal-horizon bound. The precise value
$\alpha = 0.095$ adopted in the phenomenological companion follows
from calibration of the chameleon screening profile~\cite{KhouryWeltman2004};
the formal inequality~\eqref{eq:main} requires only $\alpha > 0$, and
$\alpha$ should be treated as a prior-anchored fit parameter in any
phenomenological application.
\end{assumption}

\begin{theorem}[Main bound, sketch]
\label{thm:main}
Under Assumptions~\ref{ass:M1}--\ref{ass:M3} and the type-II setup
of Sec.~\ref{sec:setup}, Eq.~\eqref{eq:main} holds as an inequality
along the modular flow $\sigma^{R}_{\tauR}$ of $\AR$.
\end{theorem}

\paragraph{Proof sketch.}
Start from the Wall monotonicity bound on $\Sgen$, refined to the
modular (relative-entropy) form by
Faulkner--Speranza~\cite{FaulknerSperanza2024} and to the
non-semiclassical regime by Kirklin~\cite{Kirklin2025}. The
variational step $-dS_{\mathrm{rel}}/d\tauR \le |d\langle K_{R}\rangle/d\tauR|$,
following from the infimum formula for relative entropy via the
Connes-cocycle derivative (\emph{not} from Pinsker's inequality
directly), is justified on modular-flow-covariant states by the
half-sided modular pushforward of
Ceyhan--Faulkner~\cite{CeyhanFaulkner2020} together with
the Longo positivity bound~\cite{Longo2019}, and is reinforced by the
simplified QNEC proof of Hollands--Longo~\cite{HollandsLongo2025QNEC}
which bypasses analytic-continuation arguments via an explicit shape
derivative of the relative entropy. The half-sided modular
inclusion machinery underlying this step appears again in
Chandrasekaran--Flanagan~\cite{ChandrasekaranFlanagan2026} for
horizon subalgebras, where it gives a proof of the quantum focusing
conjecture in perturbative quantum gravity compatible with the
scope of our inequality. In the
type-II crossed-product factor, the RHS source admits a spectral
decomposition into a modular-commutator piece and an operator-growth
piece. M1 identifies the latter with $\kR\,\Ck[\rho_{R}]$ up to a
monotone non-negative correction, and M3 multiplies through by the
topological activator $\Theta(\PHk[\dn])$ with $\PHk$ computed on the
dequantised field from M2. The inequality is saturated when $\Ck$
reaches the Haferkamp linear-growth plateau; in that saturation limit,
Fan (2022) $\Ck \propto \exp(S_{R})$ gives the logarithmic rate
$\dot{S}_{R} \approx \dot{\Ck}/\Ck$ as a consistency check, not as a
derivation of Eq.~\eqref{eq:main}. Full derivation is deferred to the
formal companion note. $\square$

We emphasise an explicit demarcation from
Fan~(2022)~\cite{Fan2022}: their equality $\dot S_{K} = \dot C_{K}/C_{K}$
is a logarithmic \emph{equality}; our Eq.~\eqref{eq:main} is an
\emph{upper bound} which saturates to the Fan regime but does not
claim to derive the rate itself.

\begin{lemma}[Submultiplicativity of the RHS along nested regions]
\label{lem:submult}
Let $R \subseteq R'$ be two nested observer regions sharing a modular
frame, with $R' = R \sqcup (R'\setminus R)$ at the level of the
type-II crossed-product algebra (so that the CLPW nested factor
$\AR \subseteq \mathcal{A}_{R'}$ admits a faithful normal conditional
expectation $E : \mathcal{A}_{R'}\to\AR$). Writing
$\mathrm{RHS}(R) \equiv \kR\,\Ck[\rho_{R}]\,\Theta(\PHk[\dn_{R}])$ and
under Assumptions~\ref{ass:M1}--\ref{ass:M2}, the RHS of
Eq.~\eqref{eq:main} obeys the submultiplicative upper bound
\begin{equation}
 \mathrm{RHS}(R\cup R')
 \;\le\;
 \kR\,\Ck[\rho_{R}]\,\Ck[\rho_{R'\setminus R}]\,
 \min\!\Big(\Theta(\PHk[\dn_{R}]),\,\Theta(\PHk[\dn_{R'\setminus R}])\Big),
 \label{eq:submult}
\end{equation}
with equality in the $\Ck$ factor iff the reduced state factorises as
$\rho_{R'} = \rho_{R}\otimes\rho_{R'\setminus R}$ and strict inequality
in the generic entangled case.
\end{lemma}

\paragraph{Proof sketch.}
The $\Ck$ factor is controlled by the purity inequality
$\mathrm{Tr}(\rho_{R\cup R'}^{2})\ge
 \mathrm{Tr}(\rho_{R}^{2})\,\mathrm{Tr}(\rho_{R'\setminus R}^{2})$, an
elementary consequence of Cauchy--Schwarz applied to the conditional
expectation $E:\mathcal{A}_{R'}\to\AR$ in the CLPW nested
factor~\cite{CLPW2023}, with equality iff $\rho_{R'}$ factorises. In
the type-II bipartite setting this implies
$\Ck[\rho_{R\cup R'}]\le \Ck[\rho_{R}]\,\Ck[\rho_{R'\setminus R}]$ for
the 2-R\'enyi surrogate of the $k$-design / PRU complexity of
Ma--Huang~\cite{MaHuang2025}; we adopt this submultiplicativity as a
sub-postulate of M1 in the bipartite regime, since the Ma--Huang paper
does not state submultiplicativity explicitly for general type-II
factors. The $\Theta$ factor obeys a persistent-homology nesting
bound $\PHk[\dn_{R\cup R'}]\ge\max(\PHk[\dn_{R}],\PHk[\dn_{R'\setminus R}])$
(Yip \emph{et al.} 2024 nesting property on filtered simplicial
complexes~\cite{Yip2024}), so $0\le\Theta\le 1$ monotonicity gives
$\Theta(\PHk[\dn_{R\cup R'}])\le
\min(\Theta(\PHk[\dn_{R}]),\Theta(\PHk[\dn_{R'\setminus R}]))$.
Combining the two yields~\eqref{eq:submult}. The three required
properties --- submultiplicativity~(a), saturation on product states~(b),
and strict inequality on generic entangled states~(c) --- are verified
numerically on a toy bipartite system ($N_{R}=4$, $N_{R'\setminus R}=4$
qubits, XXZ modular Hamiltonian, $\beta=1$) in
\texttt{derivations/V6-lemma-submultiplicativity.py}; all three
asserts pass to $<10^{-8}$, with equality holding exactly in the
purity factor on product states and strict inequality
$\Ck[\rho_{R\cup R'}]\approx 18.2 < \Ck[\rho_{R}]\,\Ck[\rho_{R'\setminus R}]\approx 26.3$
on the entangled thermal state. $\square$

% =============================================================
\section{Dequantisation map}
\label{sec:dequant}
% =============================================================

\begin{definition}[Dequantisation map]
\label{def:dequant}
For a reduced state $\rho_{R}(\tauR)$ on $\AR$ and a smeared
number-density operator $\hat n(x)$ on $\AR$ at QRF coarse-graining
scale $\sigma_{\mathrm{cg}}$, define
\begin{equation}
 \dn(x,\tauR) \;\equiv\;
 \mathrm{Tr}_{R}\!\left[\rho_{R}(\tauR)\,\hat n(x)\right]
 \;-\; \langle n \rangle .
 \label{eq:dequant}
\end{equation}
\end{definition}

By construction $\dn$ is real (trace pairing of Hermitians), mean-zero
(subtraction of the background), and covariant under modular flow
($\sigma^{R}_{\tauR}$ acts as a one-parameter automorphism that
commutes with the partial trace in the semifinite factor). The map
$\rho_{R}\mapsto\dn$ is CPTP in its state argument. A sympy
verification of these three properties on a toy type-II$_{1}$ factor
($N=6$ qubits, $N_{R}=4$ visible, XXZ modular Hamiltonian) is reported
in Appendix~\ref{app:computation}: all three asserts pass
to $<10^{-8}$ and RMS $\dn$ at $\tauR=0.7$ is
$2.5\times 10^{-2}$.

The semi-classical recovery of $\dn$ as a Gaussian random field in the
large-$N$ QRF coarse-graining limit is postulated as M2
(Ass.~\ref{ass:M2}), following the Jacobson-2016 entanglement-
equilibrium heuristic. The persistent-homology functor $\PHk$ is then
well-defined on $\dn$ as a classical Morse-theoretic object, resolving
the category mismatch between the type-II operator algebra on which
$\rho_{R}$ lives and the classical density field on which the
Yip~et~al.\ pipeline~\cite{Yip2024} operates. Equivalently,
$\PHk[\dn]$ admits a derived-category formulation as the
Kashiwara--Schapira constructible sheaf associated to the sub-level
set filtration
$\{\dn \le t\}$~\cite{KashiwaraSchapira2018}, with
$\dim H^{k}(F_{t}) = $ the PH$_{k}$ Betti curve (KS Thm.~1.4) and
interleaving distance on $D^{b}(\mathrm{Sh}_{c}(\mathbb{R}))$ matching
the bottleneck distance on
barcodes~\cite{BerkoukGinot2018}; this home is strictly equivalent at
the Betti-number level and is occasionally preferable for formal
arguments. Moreover, in the semi-classical limit of M2, the modular
flow $\sigma^{R}_{\tauR}$ acts on $\dn$ by pullback along the
Hamiltonian isotopy $\Phi_{\tauR}$ generated by the principal symbol
of the modular Hamiltonian $K_{R}$; since this symbol is homogeneous
of degree one in the cotangent coordinate (Bisognano--Wichmann) and
the cosmological Killing vector vanishes at the de~Sitter horizon,
the isotopy satisfies the homogeneity and compact-support hypotheses
of Guillermou--Kashiwara--Schapira~\cite{GuillermouKashiwaraSchapira2012},
whose main theorem lifts $\Phi_{\tauR}$ to a unique sheaf quantisation
$Q_{\Phi_{\tauR}}$ on $D^{b}(\mathrm{Sh}_{c}(M))$. The intertwining
$F \circ \sigma^{R}_{\tauR} = Q_{\Phi_{\tauR}} \circ F$ holds as a
set-level identity on sub-level filtrations, so PH$_{k}$-covariance
under modular flow is \emph{derived under M2} rather than postulated
separately; the status of M2 itself remains ANSATZ
(Ass.~\ref{ass:M2}).\footnote{The compact-support hypothesis required
by~\cite{GuillermouKashiwaraSchapira2012} is satisfied on the dS
static patch (cosmological Killing vector vanishing at the horizon)
but not for expanding causal diamonds in generic FLRW. The
DERIVED-UNDER-M2 covariance statement is therefore strictly scoped
to the static patch; an extension to FLRW algebras of observables
(Kudler-Flam--Leutheusser--Satishchandran~\cite{KFLS2024})
requires a non-compact-support lift of GKS and is left as an open
question.} Because
Eq.~\eqref{eq:main} is an upper bound and the dequantisation map enters
only through the $\Theta$ modulation (which can only reduce the bound
from $1$ towards $0$), the exact form of $\dn$ matters only insofar as
it affects $\PHk$; an upper-envelope estimate $\Theta(\PHk) \le 1$
already preserves the bound, so no lossless reconstruction of
$\rho_{R}$ from $\dn$ is claimed or needed.

% =============================================================
\section{Relation to existing GSL statements}
\label{sec:relation}
% =============================================================

Equation~\eqref{eq:main} sits in a four-way limit diagram.

\paragraph{(L1) Wall / Faulkner--Speranza limit.}
Taking $\Theta\to 1$ recovers the Faulkner--Speranza
modular monotonicity bound~\cite{FaulknerSperanza2024}, of which our
result is a $\Theta$-weighted refinement.

\paragraph{(L2) Fan-2022 saturation and the logistic envelope.}
When $\Ck[\rho_{R}]$ approaches its Haferkamp plateau
$\Ck^{\max}\sim\exp(S_{R})$, the main bound~\eqref{eq:main} admits a
tightening that absorbs the Fan (2022) logarithmic Krylov law.

\begin{proposition}[Logistic envelope at scrambling]
\label{prop:logistic}
Under Assumptions~\ref{ass:M1}--\ref{ass:M3} and the additional
postulate that the modular-frame complexity $\Ck[\rho_{R}(\tauR)]$
evolves under a logistic-type envelope bounded by
$\Ck^{\max}\sim\exp(S_{R})$ (Brown--Susskind~\cite{BrownSusskind2018}
conjecture of complexity growth to exponential saturation,
Haferkamp~\cite{Haferkamp2022} lower bound at the linear-growth scale),
the inequality
\begin{equation}
 \frac{d\Sgen[R]}{d\tauR}
 \;\le\;
 \kR\,\Ck[\rho_{R}(\tauR)]\,
 \left(1-\frac{\Ck[\rho_{R}(\tauR)]}{\Ck^{\max}}\right)
 \,\Theta\!\left(\PHk[\dn(\tauR)]\right)
 \label{eq:logistic}
\end{equation}
holds along the modular flow of $\AR$.
\end{proposition}

\paragraph{Proof sketch.}
The RHS of~\eqref{eq:logistic} is strictly bounded above by the RHS
of~\eqref{eq:main} since the extra factor $(1-\Ck/\Ck^{\max})$ lies
in $[0,1]$; provided the source identification in M1 is refined
to a logistic rate along modular flow, any state that saturates
\eqref{eq:main} saturates~\eqref{eq:logistic} exactly at
$\Ck=\Ck^{\max}/2$, recovers Wall / Faulkner--Speranza at
$\Ck\ll\Ck^{\max}$, and reproduces the Fan (2022) asymptotic
$\dot S_{R}\to 0$ at $\Ck\to\Ck^{\max}$ where $\dot\Ck\to 0$.
Both (\ref{eq:main}) and~(\ref{eq:logistic}) are therefore consistent
with Wall and Fan; the logistic envelope is the sharper of the two
in the scrambling regime. $\square$

In the logistic regime, the Fan logarithmic rate
$\dot S_{R}\approx\dot C_{K}/C_{K}$ emerges as the leading-order
expansion of \eqref{eq:logistic} in $1-\Ck/\Ck^{\max}\ll 1$; no
competing derivation of Fan is claimed.

\paragraph{(L3) Cold-observer limit.}
For $T_{R}\to 0$ the modular flow degenerates, $\kR\to 0$, and the
bound reads $0 \le 0$: trivially consistent.

\paragraph{(L4) Classical limit / EGJ recovery.}
Taking $\dn\to 0$ so that $\PHk\to 0$ kills $\Theta$ only if
$\PHk^{\mathrm{c}}$ is held fixed; more carefully, in the strict
classical EGJ regime Eq.~\eqref{eq:main} reduces to the Clausius
form $dS = \delta Q/T$ for the geometric entropy
contribution~\cite{Jacobson1995} (Eling--Guedens--Jacobson 2006). QRF reduction: two
observers sharing a modular frame have compatible bounds by the
DEHK covariance theorem~\cite{DEHK2025a,DEHK2025b}.

\subsection{Relation to computational gravity programmes}
\label{sec:compgrav}

A parallel line of work derives gravitational dynamics from complexity
first-laws. Pedraza--Russo--Svesko--Weller-Davies~\cite{Pedraza2022Threads}
and Carrasco--Pedraza--Svesko--Weller-Davies~\cite{Carrasco2023} establish
a variational complexity/volume first-law identity that, under
stationarity of a holographic-complexity volume functional, reproduces the
linearised Einstein equations as an integrated equality. Our
Eq.~\eqref{eq:main} operates on an orthogonal axis: it is a temporal
differential inequality $d\Sgen/d\tauR \le \kR \Ck \Theta$ on the
crossed-product algebra $\AR$, formulated along the modular-flow
parameter $\tauR$ rather than as a static first-variation identity on a
bulk slice. The two statements are therefore complementary --- a
variational static route to gravity from volume-complexity
(Pedraza~et~al.) versus a differential modular-time route to an
entropy-rate bound from $k$-design complexity (this work) --- and neither
subsumes the other.

\subsection{Comparative placement}
\label{sec:comparative}

\begin{center}
\small
\begin{tabular}{lcccc}
\hline
Statement & Form & Time variable & Source term & Rigour status \\
\hline
Wall 2011~\cite{Wall2011}            & $dS/d\lambda\ge 0$     & affine $\lambda$ & $-$               & theorem (classical horizon) \\
Faulkner--Speranza 2024~\cite{FaulknerSperanza2024}
                                     & $\Delta S\ge\Delta K$   & modular cut      & relative entropy  & theorem (type II) \\
Kirklin 2025~\cite{Kirklin2025}      & refined $\Delta S$      & modular          & commutator term   & theorem (non-semicl.) \\
Fan 2022~\cite{Fan2022}              & $\dot S_{K}=\dot C_{K}/C_{K}$ & ordinary $t$   & Krylov complexity & theorem (log regime) \\
This work (Eq.~\ref{eq:main})        & $d\Sgen/d\tauR\le\kR\Ck\Theta$ & modular $\tauR$ & $k$-design $\Ck$ + PH activator & bound under M1--M3 \\
This work (Prop.~\ref{prop:logistic})
                                     & logistic envelope       & modular $\tauR$  & $\Ck(1-\Ck/\Ck^{\max})$ & bound under M1--M3 + logistic \\
\hline
\end{tabular}
\end{center}

The two rows of the present work refine the Faulkner--Speranza
modular bound by adding (i)~a complexity source identified with
$k$-design PRU~\cite{MaHuang2025}, (ii)~a topological activator
$\Theta$ anchored on persistent homology of a dequantised field, and
(iii)~a logistic envelope that absorbs the Fan logarithmic regime at
scrambling saturation. Rigour degrades monotonically down the table:
all prior rows are theorems, ours are bounds under labelled
postulates.

% =============================================================
\section{Explicit assumptions summary}
\label{sec:assumptions}
% =============================================================

\begin{center}
\begin{tabular}{lll}
\hline
Label & Status & Literature anchor \\
\hline
M1 (Ass.~\ref{ass:M1})  & POSTULATE    & Brown--Susskind 2018;
                                        Caputa--Mag\'an 2024 \\
M2 (Ass.~\ref{ass:M2})  & ANSATZ       & Jacobson 2016;
                                        \texttt{V6-dequantisation-map.py} \\
M3 (Ass.~\ref{ass:M3})  & CONJECTURAL  & Khoury--Weltman 2004~\cite{KhouryWeltman2004};
                                        Barrow 2020 \\
\hline
\end{tabular}
\end{center}

M1 identifies the modular-commutator source with $k$-design complexity;
a type-II theorem is absent from the literature. M2 is a
semi-classical CLT/LLN ansatz on QRF coarse-graining; the Gaussian
recovery is verified partially on a toy factor. M3 restricts
$\alpha \in (0, 0.1]$ as a Barrow-anchored range (the specific value
$\alpha = 0.095$ used in the phenomenological companion is a
calibration choice, not a derived quantity).

% =============================================================
\section{Programmatic outlook}
\label{sec:outlook}
% =============================================================

We close with three remarks on the conceptual context of
Eq.~\eqref{eq:main}. Each is offered as \emph{motivation} for further
work, \emph{not} as a theorem of the present paper; where an analogy or
coincidence is invoked, we state its scope and its limits.

\subsection{Entropy--complexity equivalence as a motivating analogy}
\label{sec:outlook:equivalence}

The inequality~\eqref{eq:main} juxtaposes two a priori distinct
functionals on $\AR$: the generalised entropy $\Sgen$ and the
$k$-design complexity $\Ck$. Both are monotone, trace-class on the
type-II factor, covariant under the modular flow
$\sigma^{R}_{\tauR}=\mathrm{Ad}(\Delta_{R}^{i\tauR})$, and both are
bounded by $\exp(S_R)$ at scrambling saturation. One may therefore ask
whether they play interchangeable roles in any given operational
protocol internal to $\AR$.

\begin{motivation}[Entropy--complexity indistinguishability, MOTIVATION]
\label{motiv:eq}
At the level of heuristic analogy, the observation that
Brown--Susskind~\cite{BrownSusskind2018} conjecture a second law of
complexity $\Ck\nearrow$ parallel to the GSL
$\Sgen\nearrow$~\cite{Wall2011} suggests that an operational test
internal to $\AR$ might not discriminate small variations $\delta\Sgen$
from small variations $\delta\Ck$ of the same sign. This is an
\emph{analogy} with the Einstein 1907 equivalence principle, in the
limited sense that each proposes a paired functional
$(\Sgen,\Ck)$ as the physical observable, with no separate access to
either component. We emphasise that this analogy is not a theorem: the
operational protocol that would realise the indistinguishability has
not been constructed, and the type-II semifiniteness alone does not
force $\delta\Sgen$ and $\delta\Ck$ to co-vary.
\end{motivation}

The motivating analogy frames a future programme: to construct an
explicit operational protocol, restricted to the algebra $\AR$, under
which $(\delta\Sgen,\delta\Ck)$ enters as a single measured quantity.
Success of this programme would elevate the present inequality to a
more geometric status; failure would leave it as the standalone
information-theoretic bound of Theorem~\ref{thm:main}. Neither outcome
is prejudged here.

\subsection{Modular conductance $\kR$ and its two observer-dependent scales}
\label{sec:outlook:kappa}

The Bisognano--Wichmann--Unruh identification
$\kR = 2\pi k_{B} T_{R}/\hbar$ gives a dimensional form for the
modular-flow normalisation. Two extremal values frame the range of
physically motivated observers:

\begin{itemize}[leftmargin=2em,itemsep=4pt]
 \item \textbf{Ultraviolet (Planck) observer.} With $T_{R}=T_{P}=
   \sqrt{\hbar c^{5}/(G k_{B}^{2})}$, one finds
   $\kR^{\mathrm{UV}} = 2\pi k_{B} T_{P}/\hbar = 2\pi\,\omega_{P}$
   where $\omega_{P}=\sqrt{c^{5}/(G\hbar)}\approx 1.85\times 10^{43}$
   rad\,s$^{-1}$ is the Planck angular frequency. We stress that
   $\kR^{\mathrm{UV}}$ is $2\pi$ times $\omega_{P}$, not $\omega_{P}$
   itself; the identification is dimensional, not numerical.
 \item \textbf{Infrared (de Sitter) observer.} With $T_{R}=T_{dS}=
   \hbar H_{0}/(2\pi k_{B})$ (Gibbons--Hawking), one finds
   $\kR^{\mathrm{IR}} = 2\pi k_{B} T_{dS}/\hbar = H_{0}$ exactly. This
   identification is a consequence of the definition of the de Sitter
   temperature, not an independent statement.
\end{itemize}

We do \emph{not} promote $\kR$ to the status of a fundamental
constant. Its values at the two extremal observers are determined by
$G$, $\hbar$, $c$, $k_{B}$ and the choice of observer $R$, all of
which are fixed within existing CODATA conventions. What the present
work does single out is that $\kR$ organises, within a single
formalism, quantities that were previously treated in distinct
contexts (Planck-scale thermal normalisation; de Sitter horizon
temperature).

\subsection{The ratio $(H_{0}/\omega_{P})^{2}$ as an order-of-magnitude
coincidence}
\label{sec:outlook:cosmo-ratio}

The dimensionless combination
$\left(\kR^{\mathrm{IR}}/\kR^{\mathrm{UV}}\right)^{2}
 = (H_{0}/(2\pi\omega_{P}))^{2}$
carries a numerical order that one may compare with the cosmological
vacuum ratio $\rho_{\Lambda}/\rho_{P}$, where
$\rho_{P}=c^{7}/(\hbar G^{2})$ and $\rho_{\Lambda}
= \Omega_{\Lambda} \cdot 3 H_{0}^{2} c^{2}/(8\pi G)$. With
$H_{0}=67.4$~km\,s$^{-1}$\,Mpc$^{-1}$ and $\Omega_{\Lambda}=0.6847$,
explicit numerical evaluation gives
$\log_{10}\rho_{\Lambda}/\rho_{P}\approx -122.95$ and
$\log_{10}(H_{0}/\omega_{P})^{2}\approx -121.86$; these agree at the
level of $\sim 1$ decade. We treat this as an
\emph{order-of-magnitude coincidence}, not an exact structural
identity. The gap of $\sim 1$ dex reflects the absence of a derived
prefactor relating the two ratios, and no such derivation is claimed
here. Any attempt to promote the two ratios to an exact identity
requires an explicit route from the modular ratio to the vacuum
energy density, which lies outside the scope of the present
formalism.

\subsection{Positioning with respect to related programmes}
\label{sec:outlook:programmes}

The present paper contributes one differential inequality on modular
entropy rate. It is neither intended nor advertised as a competing
unification scheme. Its overlap with adjacent programmes is:

\begin{itemize}[leftmargin=2em,itemsep=4pt]
 \item \textbf{Relational quantum mechanics and thermal
   time}~\cite{ConnesRovelli1994}. Our use of $\tauR$ as the physical
   time of observer $R$ is directly Connes--Rovelli in character; we
   add no new element to that thesis.
 \item \textbf{Jacobson--EGJ local
   Clausius}~\cite{Jacobson1995}. Our inequality is a modular-frame
   refinement of the internal-production term $d_{i}S$; we specify a
   particular functional form for the source, not a new principle.
 \item \textbf{Holographic complexity~(C$=$V,
   C$=$A)}~\cite{Pedraza2022Threads,Carrasco2023}. Orthogonal axis
   (static first-variation identity on a bulk slice) treated in
   Sec.~\ref{sec:compgrav}; see the comparative table there.
 \item \textbf{Brown--Susskind second law of
   complexity}~\cite{BrownSusskind2018}. Our M1 postulate is
   explicitly a modular-frame adaptation of this programme; the
   Caputa--Mag\'an modular-Krylov route is the natural rigorous
   upgrade.
\end{itemize}

None of these programmes is subsumed by, or subsumes, the present
work. The correct description is that our inequality refines the
modular side of the local Clausius equality of Jacobson--EGJ by
adding a complexity source with a topological activator, adapted to
the type-II crossed-product algebras of CLPW and DEHK.

% =============================================================
\section{Conclusion}
\label{sec:conclusion}
% =============================================================

We have proposed an information-theoretic upper bound on the rate of
generalised entropy along the modular flow of a type-II crossed-product
observer algebra. The bound extends the
Eling--Guedens--Jacobson~$d_{i}S$ schema into the modular regime, with
$k$-design complexity as the source and a persistent-homology
topological activator. It recovers Wall / Faulkner--Speranza
monotonicity in the $\Theta\to 1$ limit and saturates to the Fan~(2022)
logarithmic Krylov rate in the scrambling regime.

No cosmological prediction is claimed in this paper. Phenomenological
consequences, falsifier construction, and forecast analyses are
developed in the companion paper~\cite{Remondiere2026ECIv5}.
As the modular flow has a built-in direction, Eq.~\eqref{eq:main}
implies $d\Sgen/d\tauR \ge 0$ in the modular frame but does \emph{not}
purport to explain the cosmological arrow of time.

Open questions: (i)~a rigorous type-II derivation of M1 along the
Caputa--Mag\'an modular-Krylov route; (ii)~a semi-classical proof of
M2 beyond the toy-factor verification; (iii)~a first-principles
anchor for $\alpha$ in M3, most plausibly from the
Barrow fractal-horizon entropy, with the caveat that recent DESI DR2
constraints on the Barrow exponent $\Delta$~\cite{Barrow2025DESI}
find a best-fit $\Delta < 0$ while admitting $\Delta = 0$
(i.e.\ standard Bekenstein--Hawking) at $\sim 2\sigma$ — the v6
range $\alpha \in (0, 0.1]$ is therefore observationally permissible
but not empirically preferred; (iv)~extension of the KS-microlocal
covariance beyond the static patch to FLRW observer algebras along
the Kudler-Flam--Leutheusser--Satishchandran
line~\cite{KFLS2024}. Far-future structural extensions
worth exploring include recasting $\Ck$ as a Nielsen length on the
stack of $D$-modules in the Gaitsgory--Raskin~\cite{GaitsgoryRaskin2024}
proof of the geometric Langlands conjecture, and embedding the
inequality as a boundary constraint in the
Freed--Moore--Teleman~\cite{FreedMooreTeleman2024} SymTFT sandwich; a
Kapustin--Witten-type $\Ck\leftrightarrow\PHk$ duality is
structurally suggested but has no current proof beyond
$\mathcal{N}=4$ SYM on $\Sigma\times C$.

% =============================================================
\appendix

\section{Computational realisation}
\label{app:computation}
% =============================================================

The four companion scripts in \texttt{derivations/} implement the
computational checks referenced in the main text:
\begin{itemize}[leftmargin=2em,itemsep=2pt]
 \item \texttt{V6-dequantisation-map.py} --- toy type-II$_{1}$
   factor, $N=6$ qubits, $N_{R}=4$ visible, XXZ modular Hamiltonian,
   $\beta=1$. Verifies reality of $\dn$, mean-zero at $\tauR=0$, and
   commutation of the dequantisation map with modular flow to
   $<10^{-8}$.
 \item \texttt{V6-variational-derivation.py} --- sympy check of the
   variational form of the Faulkner--Speranza modular inequality in
   the type-II crossed product, under the M1 identification.
 \item \texttt{V6-krylov-to-kdesign.py} --- numerical bridge between
   Krylov complexity $C_{K}$ and $k$-design complexity $\Ck$, with
   explicit demarcation of the Fan-2022 logarithmic regime as the
   saturating limit.
 \item \texttt{V6-dimensional.py} --- sympy dimensional audit of
   Eq.~\eqref{eq:main}; confirms $[\kR]=\mathrm{nat}\cdot t^{-1}$,
   $\Ck$ dimensionless, $\PHk/\PHk^{\mathrm{c}}$ scale-free, $\alpha$
   pure number.
\end{itemize}
All sympy asserts in the four scripts pass; cross-model verification
by Gemini and Magistral derivation agents converges on
\emph{ansatz-publishable} status with the inequality form adopted here.

\section{Dimensional numerology: negative result}
\label{app:numerology}

As additional motivation for the formal-algebraic focus of the main
text, we report a brief pre-write negative result on dimensional
numerology. A systematic scan of monomial relations
$\prod_{i} (X_{i}/M_{P})^{a_{i}} \;=\; \Lambda/M_{P}^{2}$
with $|a_{i}|\le 4$ across
$X_{i}\in\{H_{0}, \Lambda^{1/4}, m_{e}, m_{p}, m_{\nu}, v_{\mathrm{EW}}, \Lambda_{\mathrm{QCD}}, F_{\pi}, \alpha\}$
produces no match below $1\%$ beyond two already-physically-motivated
anchors: the trivial Friedmann identity
$\Lambda \approx 3 H_{0}^{2}\Omega_{\Lambda}/c^{2}$, and the
dark-dimension prediction
$\Lambda^{1/4}\approx m_{\nu}\approx 2.24$~meV of
Montero--Vafa--Valenzuela~\cite{Montero2022} (within a factor
$O(1)$). The derived
proposal $\langle H\rangle \sim \Lambda^{1/6} M_{P}^{1/3}$ of
Gonzalo--Montero--Obied--Vafa~\cite{GonzaloMonteroObiedVafa2023}
fails by three orders of magnitude against the measured
$v_{\mathrm{EW}} = 246$~GeV. No exponent-rational power-law linking
$H_{0}/\omega_{P}$ to $m_{e}/M_{P}$ survives: the best fit
$k_{\mathrm{opt}} \approx 2.72$ is non-rational and physically
unmotivated. This negative result reinforces the paper's thesis that
progress at the $\Lambda$/hierarchy frontier should be sought in
formal-algebraic structure (as in Eq.~\eqref{eq:main} and
Prop.~\ref{prop:logistic}), not in coincidence-of-exponents
numerology.

% =============================================================
\bibliographystyle{apsrev4-2}
\bibliography{../eci}
% =============================================================

\end{document}
